“Dag Emmy, ik vertrek!”,
bleek eerder: “ik verrek”

samen met vluchtbeleid.
Met “je moet ook altijd”,

kon Moeder aantijgen,
zoals “je zin krijgen”,

gevolgd door “verwend nest”
als een platte scheldrest.

Emmy vroeg dan oprecht
“Waarom is dat zo slecht?”

Uiteindelijk bleek dat
Moeder graag háár zin had.

Ze zag de splinter goed,
maar er ontbrak de moed

om de balk in haar oog
te zien in haar betoog.

Dat gold niet voor Emmy,
die zag de projectie.

Moeder mocht graag preken.
Uren kon ze spreken.

Viel in herhalingen
over bepalingen,

die een profeet haar gaf.
Emmy vroeg zich soms af:

“Mama, ik heb godswoord
nu al zo vaak gehoord

dat ik er gek van word”.
Moeder snauwde dan kort:

“en wat denk je van mij!
Ik word ook niet echt blij

dat je het maar niet snapt
wat ik je heb verklapt

hoe vaak ik het ook zeg.”
Emmy vond niet haar weg

in een geleerd dogma,
maar zocht anubhava:

de ervaren waarheid
dat leven echt belijdt.

Zij had als strategie:
“Geen protest, let it be”,

geleerd van haar ouders.
Zij van haar voorouders.

Ze dacht ook aan haar oom.
“Zelfs een stokoude boom

schiet bloot groene blaadjes.
Een bij helpt bij zaadjes

bij puur jonge bloemen”,
wist hij wulps te zoemen.

Emmy, gehersenspoeld,
peinsde zoals bedoeld:

“Verzet is toch zinloos.
Het maakt de jager boos.

Ik geef me dus over.
Is hij een struikrover?

Dan word ik vast beroofd!
Hoewel hij niet gelooft,

dat ik maar iets bezit.
Anders dan een proefrit.

Ik zal mijn lijf geven
om weer vrij te leven.

Lot bood zijn dochters aan.
De Leviet mocht vrij gaan

door voor een paar uren
zijn bijvrouw te sturen

onder de vleugelen,
om te beteugelen

de woedende massa.
Het lijf als activa.”

Moeder duwde pa aan.
Hij snurkte weer voldaan

na zijn daagse paring.
Daar het niet overging,

moest hij uit de bedstee.
Hij luisterde gedwee.

Maar ook zonder de schurk
hoorde ze weer gesnurk.

Ze draaide zich boos om
op zoek naar het gebrom,

maar trof een leeg bed aan.
Ze liet het voorbijgaan

en zakte weer in slaap.
Haar aandrang tot een gaap

voelde als gerochel
en klonk als een orgel.

“L’enfer, c’est les autres”,
zei ooit Jean-Paul Sartre.

Dat mannen ontladen
van dreigende kwaden

door zich te vergrijpen,
kon ze niet begrijpen.

Tot haar te binnen schoot:
Ze werd ooit deelgenoot

van een missionaris
dat er zelfs een aap is

die elk conflict oplost
of een soort schuld inlost

door andere lieden
een wip aan te bieden

tot er kalmte ontstaat.
Dit dempt de zin in kwaad.

De jager behagen
door zich te gedragen

als een welwillend dier
voor een beetje plezier

kon Emmy ook helpen
haar angstgolf te stelpen.

Vrijheid en veiligheid
binden tot vrijigheid.

“Was Moeder als dienstmeid
niet ook tot veel bereid?

Mannen waren vast erg
in die louche herberg!”

Al was moeder moe der
preken, bleef de moed er

toch onvermoeibaar in.
Dus sprak ze weer een zin

dit keer zelfs in haar slaap.
Het snurken in een gaap

verried haar zenuwtrek.
Dus vond ze het niet gek

toen vader haar aansprak
dat ze in haar slaap sprak.

Biechten dat ze snurkte
kon de steeds gehurkte,

die al vaak moest bukken,
vooralsnog niet lukken.

Toen moeder wat later
ronkte uit haar snater

bedacht vader zijn wraak
gekoppeld aan vermaak.

Hij schoof bovenop haar
en maakte zijn lid klaar

om met één harde stoot
te landen in haar schoot.

Moeder schrok suf wakker
en de brave stakker

vroeg stomverbaasd: “waarom?”.
“Hé? Je vroeg er zelf om!”

Moeder van Zelf sprekend
zei wat het betekent

de uitspraak die Jezus ,
althans naar Matteüs,

deed over kinderen
en het verhinderen
van elke toegang tot
het koninkrijk van God.

Moeders zelfhulpstudie,
ofwel theologie,

gaf bij Emmy’s protest
toen Lucas als attest.

Ze stelde; “stelligheid
is jouw onzekerheid

dat zichzelf niet wil zien.
Het is lomp bovendien!”

Ze schopte uit woede
Emmy’s achterhoede.

Uit ervaring geleerd
had Emmy zich verweerd

door een houten plankje
als pijnlijk bedankje

achter in haar rokje
te doen als stootblokje.

Het was misschien gemeen
maar de gebroken teen

weerhield moeder een tijd
van wrede geloofsstrijd.

Moeder had sensaties
zonder associaties.

Zonder het label bang.
Dat was voor Emmy wrang,

want Moeder dumpte het
in elk lelijk facet

op Emmy haar bordje.
Moeder haar gifstortje

was altijd religie.
Ze bekeerde Emmy

in plaats van te voelen.
In plaats van te woelen

in de gewaarwording.
Het werd geen eenwording

door een observatie
van elke sensatie.

Er was nog een conflict.
Het brein die het niet pikt.

Het dualisme hield
daar het aandacht onthield.

Er was wel geen verhaal.
Geen afstand door verbaal

toekomst en verleden
aan het hart te smeden.

Maar ook geen verbinding
door haar bevrediging

middels het botvieren
als devoot versieren.

Geen kuise versmelting,
maar juist verwijdering.

Reflectieve afstand
kwam wel van Emmy’s kant.

Ze stelde soms vragen
om haar uit te dagen.

Moeder gaf nooit antwoord,
maar voortdurend Godswoord,
hoe ze ook ageerde.
Erich Fromm typeerde

ooit dit dogmatisme;
‘afweermechanisme’.

“Wat wordt de reactie
op mijn vraag?”, vroeg Emmy

zich afgezonderd af
toen de Jager zijn staf

zijn broekspleet liet bollen
en haar bloed deed stollen.
